A arquitetura \gls{crnn} é amplamente utilizada em aplicações que exigem 
inferência de classes em sequência, 
pois combina um extrator convolucional (\gls{cnn}) 
a fim de obter representações visuais, 
uma camada recorrente (\gls{rnn}), com o objetivo de modelar dependências contextuais ao longo do eixo horizontal, 
e a função de perda \gls{ctc}, com a finalidade de alinhar as previsões da rede à sequência de rótulos sem a necessidade de 
segmentação explícita dos caracteres \cite{shi2017crnn}.

A técnica \gls{ctc} resolve o problema central de sequências de comprimento variável: 
em geral, não se sabe onde cada caractere começa e termina na imagem, 
e o número de \textit{frames} produzidos pela CNN é maior do que o número de caracteres 
da palavra ou número a ser reconhecido \cite{graves2006ctc}. 
Para lidar com isso, a técnica CTC introduz dois mecanismos fundamentais: 
um \textit{token} oculto de espaçamento (\textit{blank}) e uma 
função de colapso que mapeia a sequência de previsões da rede 
para o rótulo final \cite{graves2006ctc, hannun2017ctc}.

\subparagraph{O \textit{token} \textit{blank} e o colapso de repetições.}
A cada passo de tempo, 
a rede produz uma distribuição de probabilidade sobre o alfabeto 
acrescido do símbolo especial \textit{blank} (\(\epsilon\)). 
O \textit{blank} representa a ausência de emissão de caractere naquele instante, 
funcionando como um separador entre caracteres idênticos adjacentes 
\cite{graves2006ctc}. 
Após a rede produzir a sequência completa de previsões, 
aplica-se uma função de colapso \(\mathcal{B}\) 
que (i) remove caracteres repetidos consecutivos e 
(ii) remove os \textit{blanks} \cite{hannun2017ctc}. 
Por exemplo, a sequência de caminho \(--h--e--l--l--o--\) colapsa 
para \(\text{``hello''}\); 
sem o \textit{blank}, as repetições de ``l'' seriam indistinguíveis de 
um único ``l'' \cite{hannun2017ctc}.

\subparagraph{Marginalização de alinhamentos.}
O treinamento da CTC não exige um alinhamento predefinido entre 
os \textit{frames} e os caracteres. Em vez disso, 
a função de perda calcula a probabilidade da sequência de rótulos correta como 
a soma das probabilidades de todos os caminhos de alinhamento válidos que, após o colapso, 
produzem aquele rótulo \cite{graves2006ctc}. 
Essa marginalização é feita de forma eficiente pelo algoritmo \textit{forward-backward}, 
que explora a estrutura do problema para evitar a 
enumeração explícita de todos os caminhos, 
que seria computacionalmente inviável \cite{graves2006ctc, hannun2017ctc}. 
Assim, a rede pode ser treinada diretamente com pares (imagem, transcrição), 
sem anotações de posição caractere a caractere.

\subparagraph{Inferência.}
No momento da inferência, a rede produz, para cada passo de tempo, 
a classe mais provável (ou uma distribuição sobre as classes). 
A sequência final é obtida pelo colapso das previsões, 
opcionalmente combinado com busca em feixe (\textit{beam search}) 
para explorar múltiplos alinhamentos e melhorar a acurácia \cite{hannun2017ctc}. 
O artigo de Hannun (2017) apresenta uma explicação visual detalhada desse processo, 
incluindo exemplos passo a passo de como o \textit{blank} e o 
colapso atuam em sequências reais \cite{hannun2017ctc}.

\subparagraph{Por que a CTC é adequada ao OCR.}
Em OCR de palavras, números de medidores ou placas, 
o número de caracteres varia e não há garantia de que 
cada caractere ocupe exatamente um número fixo de \textit{frames} 
na saída da CNN. 
A CTC elimina a necessidade de segmentação manual e permite que a 
CRNN seja treinada de ponta a ponta com sequências de comprimentos diferentes, 
mantendo a coerência entre as previsões vizinhas e evitando 
o colapso de repetições idênticas \cite{shi2017crnn, graves2006ctc}.